\scanpage{116}
nb de triangles épinglés $=$ nb d'arcs

\[ \mathrm{trep}.=6(1+2c)(1+c)=6(1+3c+2c^2) \]

Nb de bitriangles épinglés : \deleted{\unsure{bitri}}

\[
\begin{aligned}
\mathrm{bitrep}&=2c^2\,\mathrm{trep}=12c^2(1+2c)(1+c)\\
&=12c^2(1+3c+2c^2)
\end{aligned}
\]

\uline{NB} Si $c=2$, \deleted{\unsure{les}} autom de $S=(S,A)$ qui \deleted{\unsure{est}} \deleted{\unsure{induit}} \unsure{invarie|invari|induisent} un bitriangle \uline{épinglé} est l'identité. Donc le grpe des autom $G$\deleted{\illegible{1}} (et $=$ le grpe \deleted{\unsure{le}} des autom \lq\lq intérieurs\rq\rq{} i.e. le ss-grpe de $G$ engendré par les $\sigma_{\underline t}$, $\underline t\in T$) est \uline{simplement transitif} sur l'ens des bitriangles, épinglés. Donc ici

\[
\begin{aligned}
G=G^{\unsure{0|\circ}},\quad\text{d'ordre}\quad &12c^2(1+2c)(1+c)\\
&=12.\deleted{\unsure{2}}4.5.3\\
&=2^4.3^2.5=720
\end{aligned}
\]

\uline{NB} Si $c=1$, $G^0\simeq\mathfrak S_3\times\mathfrak S_{\unsure{5|3}}$ est d'indice $2$ dans

\[
\begin{aligned}
G&\simeq\mathfrak S_2.(\mathfrak S_3\times\mathfrak S_3),\quad\text{qui est d'ordre}\\
&72=2^3\,3^2
\end{aligned}
\]

On l'obtient aussi comme le nb des bitriangles épinglés (puisque $G$ est \unsure{simple|simple\textsuperscript{t}|simpl.} transitif dessus) qui est

\[ 12c^2(1+2c)(\unsure{1|R}+c)=12.1.3.2=72 \]

$\Downarrow$

Si $c\geq1$, le nb de bitriangles $=\frac1{72}$ nb de bitriangles épinglés

\[ =\tfrac16\,c^2(1+2c)(1+c) \]

\margin{$\begin{array}{r}\unsure{46|16}\\\unsure{45}\\\hline\unsure{640}\end{array}\qquad720$}

% archive sheet number: 115
